\section{The question and the analytical structure}\label{sec:framework}

The purpose of the boundary bootstrap is to improve a known enclosure
of the zeros by observing the field just outside a proposed smaller
enclosure. Two kinds of information can support such a step. A uniform
relative approximation to a nonvanishing carrier gives direct zero
exclusion. A smaller amount of boundary information can suffice through
Green's formula, provided that the value at the observation point has an
independent lower bound. The second route keeps cancellations that the
first route may discard.

This manuscript develops the results of the height and time dependent
boundary investigation through its E11 review. Its starting analytic
input is the corrected complex Riemann--Siegel estimate established in
the repository's heat-density investigation. We state the input precisely
in Section~\ref{sec:expansion}, retain its review status, and derive the
positive-time extension needed here. The other imported inputs are the
classical logarithmic-phase exponent-pair estimate and two published
heat-flow zero theorems. The resulting deductions concern the complete
canonical field, not a truncated theta model.

\subsection{The field and the three time conventions}

For real $t$ and complex $z$, set
\begin{equation}\label{eq:canonical}
 H_t(z)=\int_0^\infty \ee^{tu^2}\Phi(u)\cos(zu)\dd u,
 \qquad
 \Phi(u)=\ee^u\sum_{n\ge1}
 (2\pi^2n^4\ee^{8u}-3\pi n^2\ee^{4u})\ee^{-\pi n^2\ee^{4u}}.
\end{equation}
The normalization, derived below, is
\begin{equation}\label{eq:normalization}
 H_0(z)=\frac18\xi\left(\frac{1-\ii z}{2}\right),\qquad
 \partial_tH_t=-\partial_z^2H_t.
\end{equation}
Thus $|x|/2$ is the ordinary zeta ordinate when $z=x+\ii y$.
Throughout, $\log$ is real on positive arguments and $\Log$ is the
branch on the right half-plane continued from the positive axis.
We use
\[
 L_X=\log\frac{X}{4\pi},\qquad t_X=L_X^{-1},\qquad
 \tau_X=\frac X2+\frac{\pi}{8L_X}.
\]
Three different assertions must be distinguished. In a spatial contour
at center $X$, the field is $H_{t_X}$ at one fixed time. In a center
average, $X$ varies and the family follows the curve $t_X$. In a
common-time exterior theorem, one fixed $t$ is used at every height.
The function $z\mapsto H_{t_{\Re z}}(z)$ is not used as a holomorphic
field in a contour identity.

\subsection{Main conclusions and the unresolved premise}

The strongest acquired enclosure is the following quantitative tail
statement. The word ``acquired'' means proved from the named inputs in
this manuscript and its source notes, with the internal reviews described
in Section~\ref{sec:evidence}.

\begin{theorem}[Exterior enclosure]\label{thm:main-width}
There is a finite, unevaluated $t_*>0$ such that
\[
 0<t\le t_*,\quad |x|\ge4\pi\ee^{1/t},\quad H_t(x+\ii y)=0
 \quad\Longrightarrow\quad |y|<\frac{567}{1000}.
\]
There is also a finite onset $X_*$ such that
$H_{t_X}(X+\ii y)=0$, $X\ge X_*$, implies the same strict width.
\end{theorem}

The proof in Section~\ref{sec:width} yields a relative error at most
$1/64$ on its full comparison rectangles. The width improves the
earlier $4/7$ by $31/7000$, at its own onset. It does not specify the
true smallest onset, and it gives no exclusion in the omitted compact
region.

For the Green route, define
\begin{equation}\label{eq:raw-tail-intro}
 \begin{split}
 T_X(v)&=\sum_{\lfloor\ee^{L_X/32}\rfloor<n\le
                    \lfloor\ee^{L_X/2}\rfloor}
 w_n(X)\ee^{\ii(\tau_X+v/2)\log n},\\
 w_n(X)&=n^{-65/64}\exp\left(\frac{(\log n)^2}{4L_X}\right),\qquad
 Q(X)=\int_\R P(v)|T_X(v)|^2\dd v,
 \end{split}
\end{equation}
where the positive strip kernel $P$ is specified in
Section~\ref{sec:moment} and has mass $2/5$.

\begin{proposition}[Conditional deeper step]\label{prop:intro-green}
If $Q(X)\le4/66^2$ for every sufficiently large real $X$, then every
zero $X+\ii y$ of $H_{t_X}$ at sufficiently large $X$ satisfies
$|y|<9/16$.
\end{proposition}

The proposition is a sufficient arithmetic reduction. Its premise
remains unproved, and no converse is asserted. The unconditional
moment estimates obtained below are
\[
 \|T_X-M_X\|_P\ll X^{-1/8192},\qquad
 \frac1U\int_U^{2U}Q(X)\dd X\ll U^{-25/1024},
\]
where $M_X$ is the middle band $3/64<\log n/L_X\le5/16$.
For every interval $I\subset[U,2U]$ of length $h\ge1$,
\[
 \frac1h\int_I Q(X)\dd X
 \ll U^{-1/4096}+U^{-21/512}+\frac{U^{5/128}\log U}{h}.
\]
These estimates become small on specified growing windows. They do not
bound the moment at every center. Section~\ref{sec:means} proves the
stronger, but still averaged, statement for positive resonator weights
of integer support $R\le U^\delta$, for each fixed $0<\delta<123/128$.

At the shorter scale $h_U=U^{1/32}$, define the probability measure
\[
 \dd\nu_{h_U}(s)=h_U^{-1}(1-|s|/h_U)_+\dd s,
 \qquad A_U(Y)=\int Q(Y+s)\dd\nu_{h_U}(s).
\]
An actual witness $A_U(Y)\ge5/66^2$ along an unbounded sequence
would contradict the every-center sufficient budget. It has not been
constructed or excluded. Section~\ref{sec:packets} proves instead that
a specified sparse family of aligned logarithmic packets cannot change
the complete moment by a fixed positive amount. Its explicit quartic
example achieves actual phase alignment but has insufficient mass.

\subsection{Scope and notation}

All implicit constants are uniform in the real centers and in the
fractional parts of the indicated cutoff values. They may depend on
declared fixed strip widths, a fixed finite clock window, a specified
exponent pair, or a fixed parameter written in the subscript.
An onset is chosen after these constants are fixed. An assertion valid
for each fixed width does not imply a uniform onset as that width shrinks.

The classification used here is a quantitative positive-time enclosure,
an explicitly conditional further step, and auxiliary estimates that
exclude named estimating classes. None is counted as a proof of RH or
as an iteration to the real axis. The repository result and avoidance
indices~\cite{RepoResults,RepoAvoid} route respectively to the inherited
claims and the relevant failure patterns. The specific control screen
and the scope of every negative conclusion appear in
Section~\ref{sec:evidence}; this follows the criteria in patterns
01, 02, 07 and 12~\cite{RepoAvoidPatterns}.
